\documentclass[10pt,a4paper]{amsart}
\usepackage[margin=19mm]{geometry}
\usepackage{amsmath,amssymb,mathtools}
\usepackage[scr=rsfs,cal=euler]{mathalfa}
\usepackage{xfrac}
\usepackage[unicode,hidelinks,bookmarks=false]{hyperref}
\newcommand{\ie}{i.e.\ }
\newcommand{\cf}{cf.\ }
\newcommand{\resp}{respectively\ }
\newcommand{\Subsection}{\subsection}
\providecommand{\nobreakdash}{\nobreak\hbox{-}}
\setlength{\emergencystretch}{2em}
\multlinegap0pt
\usepackage{amssymb}
\usepackage{xcolor}
\usepackage{float}
\usepackage{enumerate}
\usepackage{mathtools}
\usepackage{tcolorbox}
\usepackage{tikz}
\newtheorem{proposition}{Proposition}
\newtheorem{theorem}{Theorem}
\newcommand\C{\ensuremath{\mathbb{C}}}
\newcommand\N{\ensuremath{\mathbb{N}}}
\newcommand\Oph{\ensuremath{\mathrm{Op}_h}}
\newcommand\R{\ensuremath{\mathbb{R}}}
\title{Observability and Semiclassical Control for Schrödinger Equations on Non-compact Hyperbolic Surfaces}
\author{Xin Fu, Yulin Gong, Yunlei Wang}
\date{Source: arXiv:2602.14316v2 (CC BY 4.0)}
\begin{document}
\maketitle

\noindent\emph{Source and licence.} Observability and Semiclassical Control for Schrödinger Equations on Non-compact Hyperbolic Surfaces. Version arXiv:2602.14316v2 (CC BY 4.0), distributed by arXiv under the Creative Commons Attribution 4.0 International licence, http://creativecommons.org/licenses/by/4.0/, as recorded in the arXiv licence metadata for this exact version and shown on the abstract page https://arxiv.org/abs/2602.14316v2. Full licence notice: \texttt{licenses/arxiv-2602.14316-CC-BY-4.0.txt}.\par\smallskip

\noindent\textit{Editorial scope.} Counted unit: the article's theorem mainthmofquantized (source0.tex lines 2053--2059). The article's complete original proof of it is delivered with it (source0.tex lines 2060--2163, one \texttt{proof} environment of 104 lines). No mathematics is written, repaired or re-derived: the statement and the proof are the article's own lines, in the article's own notation, and no retained line is substituted or re-flowed.  Retained uncounted as the source-local context the proof needs: lines 1596--1598; lines 1600--1602; lines 1727--1732; lines 1881--1883; lines 1895--1903; lines 1958--1960; lines 2003--2005; lines 2023--2025.  Nothing was omitted from any retained span, so the package carries no omission record.  No retained line contains a citation, so the wrapper carries no bibliography and no bibliography entry.  No retained line contains a figure environment, an image-inclusion command or drawing code, so no image is delivered and the package declares no asset dependency.  Editorial formatting only, in addition: an AMS \texttt{amsart} preamble that restates the article's own declarations and macros used by the retained lines, namely the environments \texttt{proposition}, \texttt{theorem} on separate counters, together with the standard packages those lines need.  The retained environments print in this document as Proposition 1, Theorem 1 in order of appearance, whereas the article prints the same items with the numbering of its own full document.  The complete native source of the article as distributed in the arXiv e-print for this exact version is bundled unchanged as \texttt{sources/194702/source0.tex}.
\begin{equation}\label{pseudooflocal}
(\varphi^{-1})^*\Phi \chi A (\mathcal{H};\rho;h) \chi \Phi^{-1} \varphi^*=\Oph(a_{\varphi,\chi})\otimes \mathrm{Id}_{\mathcal{H}}.
\end{equation}

\begin{equation}\label{disresi}
    \psi A \psi^{\prime}=\mathcal{O} (h^{\infty} )_{\Psi^{-\infty,sc}(M;F^\rho)}.
\end{equation}

\begin{proposition}\label{uniformboundedtheorem}
 For any $a(x, \xi ; h) \in S^k\left(T^* M\right)$ and any $s \in \mathbb{R}$, there exists a constant $C=C(a, s, k)>0$, independent of $( \mathcal{H},\rho) \in \mathcal{C}$ and $0<h\leq 1$, such that
\begin{equation}
\left\|\mathrm{Op}_h^\rho(a)\right\|_{H_h^s\left(M;F^{\rho}\right) \rightarrow H_h^{s-k}\left(M;F^{\rho}\right)} \leq C .
\end{equation}
\end{proposition}

\begin{equation}\label{ellipticparametrix}
A=PQ+\mathcal{O}(h^{\infty})_{\Psi^{-\infty,sc}(M;F^\rho)}=Q^{\prime}P+\mathcal{O}(h^{\infty})_{\Psi^{-\infty,sc}(M;F^\rho)},
\end{equation}

\begin{proposition}\label{inverseofelliptic}
Let $k\in \R$ and $P \in \Psi_{h}^{k,sc}(M;F^\rho)$. By Proposition \ref{uniformboundedtheorem}, for any $s\in \R$,
\begin{equation*}
    P(\mathcal{H};\rho;h):H^{s+k}_h(M;F^\rho) \rightarrow H^{s}_h(M;F^\rho) 
\end{equation*}
is bounded uniformly in $(\mathcal{H},\rho,h)$. If $P$ is elliptic, then there exists $h_0\in (0,1]$, independent of $(\mathcal{H},\rho)$, such that for any $h\in (0,h_0)$, 
$P(\mathcal{H};\rho;h)$ is invertible and the inverse
 $P^{-1} \in \Psi_{h}^{-k,sc}(M;F^\rho)$.
\end{proposition}

\begin{equation*}\label{linearsymbol}
L_{\ell(x,hD_{x}),\widetilde{\chi}}:=\varphi^*\Phi^{-1}\widetilde{\chi} \big(\ell(x,hD_{x})\otimes \mathrm{Id}_{\mathcal{H}} \big) \widetilde{\chi}\Phi(\varphi^{-1})^* .
\end{equation*}

\begin{equation}\label{propofad}
\mathrm{ad}_{L}P^{-1}=P^{-1}(\mathrm{ad}_{L}P)P^{-1},\quad \mathrm{ad}_{L}(AB)=(\mathrm{ad}_{L}A)B+A(\mathrm{ad}_{L}B).
\end{equation}

\begin{equation}\label{HSformula}
f(P)= \frac{1}{\pi } \int_{\mathbb{C}} \overline{\partial} \widetilde{f}(z)(P-z)^{-1} \,L(d z),
\end{equation}

\begin{theorem}\label{mainthmofquantized}
For any $k \in \mathbb{N}$, let $P\in \Psi^{k,sc}_h(M;F^\rho)$ be a self-adjoint operator on $L^2(M;F^{\rho})$ such that $P+i$ is elliptic in $\Psi^{k,sc}_h(M;F^\rho)$ and $\sigma_{h}^{\rho}(P)$ is independent of $h$. Then, for any $f\in C_{0}^{\infty}(\R)$, 
\begin{equation}
f(P) \in \Psi^{0,sc}_{h}(M;F^\rho) \quad \mathrm{with} \ \sigma_{h}^{\rho}(f(P))=f(\sigma_{h}^{\rho}(P)),
\end{equation}
Moreover, $f(P)$ is compactly microlocalized, and $\mathrm{WF}_{h}(f(P))\subset \{(x,\xi)\in T^*M: \sigma^\rho_{h}(P) \in \mathrm{supp}\,f\}$.
\end{theorem}

\begin{proof}
For simplicity of notations, below we write $H^s_h$ for $H^s_h(M;F^\rho)$ and $A\lesssim B$ for $A\leq CB$ for some constant $C>0$ independent of $(\mathcal{H},\rho,h,z)$. In this proof, we always assume that $\mathrm{Im} \, z \neq 0$. We divide the proof into several steps.

\emph{Step 1. We show that}
\begin{equation}\label{sobolevnormofresolvent}
\|(P-z)^{-1}\|_{H^{n}_h\to H^{n+k}_{h}}\lesssim |\mathrm{Im} \, z|^{-2^n} \langle z\rangle^{2^n}
 , \qquad \forall n\in \N.
\end{equation}

Since $P \in \Psi^{k,sc}_{h}(M;F^\rho)$ is self-adjoint on $L^2(M;F^{\rho})$, by the spectral theorem, we have 
\begin{equation}\label{spectraltheorem}
\|(P-z)^{-1}\|_{L^2\to L^2} \lesssim |\mathrm{Im} \, z|^{-1}.
\end{equation}
Since $P+i$ is elliptic in $\Psi^{k,sc}_h(M;F^\rho)$, $(P+i)^{-1}: L^2(M;F^{\rho})\to H^{k}_h(M;F^{\rho})$ is uniformly bounded in $(\mathcal{H},\rho,h)$. For any $u\in C^{\infty}(M;F^{\rho})$, applying \eqref{spectraltheorem} yields
\begin{equation*}
\|u\|_{H^{k}_h}\lesssim \|(P+i)u\|_{L^2}\lesssim \|(P-z)u\|_{L^2}+\langle z\rangle\|u\|_{L^2}\lesssim |\mathrm{Im} \, z|^{-1} \langle z\rangle \|(P-z)u\|_{L^2}.
\end{equation*}
This gives \eqref{sobolevnormofresolvent} for $n=0$. We proceed by induction. Assume that \eqref{sobolevnormofresolvent} holds for $n\in \N$. Let $D=\Oph^{\rho}(\langle \xi \rangle) \in \Psi^{1,sc}_h(M;F^\rho)$. Then,
\begin{equation*}
    D(P-z)^{-1}=(P-z)^{-1}D+(P-z)^{-1}[P,D](P-z)^{-1},
\end{equation*}
where $[P,D]\in h\Psi^k_h(M;F^\rho)$. For any $u \in H^{n+1}_{h}$, we have
\begin{equation*}
\begin{aligned}
\|(P-z)^{-1}u\|_{H^{n+k+1}_h} &\lesssim \|D(P-z)^{-1}u\|_{H^{n+k}_h}\\
& \leq\|(P-z)^{-1}Du\|_{H^{n+k}_h}+\|(P-z)^{-1}\|_{H^{n}_h\to H^{n+k}_{h}}  \|[P,D](P-z)^{-1}u\|_{H^{n}_h}\\
&\lesssim  \frac{\langle z\rangle^{2^n}}{|\mathrm{Im} \, z|^{2^n}}\|u\|_{H^{n+1}_h}+\frac{\langle z\rangle^{2^n}}{|\mathrm{Im} \, z|^{2^n}} \cdot \|[P,D]\|_{H^{n+k}_h\to H^n_h} \cdot\frac{\langle z\rangle^{2^n}}{|\mathrm{Im} \, z|^{2^n}}\|u\|_{H^n_h}\\
&\lesssim |\mathrm{Im} \, z|^{-2^{n+1}} \langle z\rangle^{2^{n+1}} \|u\|_{H^{n+1}_h}.
\end{aligned}
\end{equation*}
By induction, we obtain \eqref{sobolevnormofresolvent}.

\emph{Step 2. We show that $f(P)\in \Psi^{0,sc}_h(M;F^\rho)$.} 

We first verify \eqref{disresi} for $f(P)$. Let $\psi, \psi^{\prime} \in C^{\infty}(M)$ such that $\operatorname{supp} \psi \cap \operatorname{supp} \psi^{\prime}=\varnothing$. Let 
\begin{equation*}
    p_{0}=\sigma^{\rho}_{h}(P).
\end{equation*}
It follows that $|p_{0}(x,\xi)|\geq c\langle \xi \rangle_{x}^k-1$ for some $c>0$, and hence $P -z \in \Psi^{k,sc}_{h}(M;F^\rho)$ is elliptic. By Proposition \ref{inverseofelliptic}, $(P-z)^{-1} \in \Psi^{-k,sc}_{h}(M;F^\rho)$. In particular, $\psi (P-z)^{-1} \psi'$ is locally fiberwise diagonal. Using the Helffer-Sj\"{o}strand formula \eqref{HSformula} and the estimate \eqref{spectraltheorem}, we deduce that $f(P)\in \mathcal{L}(L^2(M;F^{\rho}))$ and $\psi f(P)\psi'$ is locally fiberwise diagonal. 

By the existence of elliptic parametrix \eqref{ellipticparametrix}, for any $T\in \mathbb{N}$, there exists $Q_{T}(z) \in \Psi^{-k,sc}_{h}(M;F^\rho)$ and $R_{T+1}(z) \in \Psi^{-T-1,sc}_h(M;F^\rho)$ such that
\begin{equation*}
    I=(P-z)Q_{T}(z)+h^{T+1}R_{T+1}(z) .
\end{equation*}
For $T=0$, $Q_{0}(z)=\Oph^{\rho}((p_{0}-z)^{-1})$. Computing the symbol of $R_{T+1}(z)$ under cutoff normal charts yields that, for any $T\in \N$ and any $s\in \R$, there exists $K=K_{s,T}>0$ such that
\begin{equation}\label{sobolevnormQR}
\|Q_{T}(z)\|_{H^{s}_{h}\to H^{s+k}_{h}}  + \|R_{T+1}(z)\|_{H_{h}^{s}\to H^{s+T+1}_{h}} \lesssim |\mathrm{Im}\, z|^{-K}.
\end{equation}
For any $N\in \N$, taking $T\geq 2N+|k|$, there exists $K=K_{N,T}>0$ such that
\begin{equation}
\begin{aligned}
&\|\psi  (P-z)^{-1}\psi'\|_{H^{-N}\to H^{N}}\\
%&\leq \|\psi  Q_{T}(z)\psi'\|_{H^{-N}\to H^{N}}+h^{T+1}\|\psi (z-P)^{-1}R_{T+1}(z)\psi'\|_{H^{-N}\to H^{N}}\\
& \lesssim \|\psi  Q_{T}(z)\psi' \|_{H^{-N}\to H^{N}}+h^{T+1}\|(P-z)^{-1}\|_{H_{h}^{-N+T+1}\to  H_{h}^{N}}\cdot\|R_{T+1}(z)\|_{H^{-N}_h\to H^{-N+T+1}_h} \\
&\lesssim  h^{T}|\mathrm{Im}\, z|^{-K }\langle z\rangle^{K },
\end{aligned}
\end{equation}
where we used \eqref{sobolevnormofresolvent} and \eqref{sobolevnormQR}. By Helffer-Sjöstrand formula \eqref{HSformula}, we obtain that there exists a constant $C=C_N>0$ such that
\begin{equation}\label{cutoffoffp}
\|\psi f(P) \psi'\|_{H^{-N}_h\to H^N_h}\leq C_N h^N.
\end{equation}
Therefore, \eqref{disresi} holds for $f(P)$.

We now verify \eqref{pseudooflocal} for $f(P)$. For any cutoff normal chart $(\Phi, U,\varphi,\chi)$, let $\ell(x,hD_x),\, \widetilde{\chi}$, $\chi'$ and $ L_k$, $k=1,\cdots,N+N'$, be defined as in the proof of Proposition \ref{inverseofelliptic}. By \eqref{propofad} and \eqref{sobolevnormofresolvent}, for any $s\in \R$, there exists $K=K_{s,k}>0$ such that
\begin{equation}\label{adofresolvent}
 \big\|\operatorname{ad}_{L_{\ell(x,hD_x),\widetilde{\chi}}}(P-z)^{-1} \big\|_{H^{s}_{h}\to H^{s+k}_{h}} \lesssim h|\mathrm{Im}\, z|^{-K}\langle z \rangle^{K}.
\end{equation}
Combining \eqref{sobolevnormofresolvent} and \eqref{adofresolvent}, we find a constant $K=K_{N,N',k}>0$ such that
\begin{equation}\label{adestiP-z}
\big\|\operatorname{ad}_{L_{1}} \cdots \operatorname{ad}_{L_{N+N'}} \chi (P-z)^{-1}\chi \big\|_{L^2\left(M;F^\rho\right) \rightarrow H_h^{k+N}\left(M;F^{\rho}\right)} \lesssim h^{N+N'}|\mathrm{Im}\, z|^{-K}\langle z\rangle^{K} .
\end{equation}

%Next we apply the semiclassical Beals Theorem to show $f(P) \in \Psi^{0,sc}_{h}(M;F^\rho)$. For any linear symbol $\ell(x,hD_x)\in \Psi^{1}_h(\R^d)$, $L_{\ell(x,hD_x),\widetilde{\psi}}\in \Psi^{1,sc}_h(M;F^\rho)$ is defined in \eqref{linearsymbol}. For any $N>0$, we take the same cutoff functions $\psi_{i} \in C_{0}^{\infty}(\varphi(U))$ for $i=1,\cdots,N$ such that $\psi_{i}\equiv 1$ on $\mathrm{\psi_{i+1}}$ and $\psi_{N+N^{\prime}}\equiv 1$ on $\mathrm{supp}((\varphi^{-1})^*\psi)$, and $\psi^{\prime}\in C_{0}^{\infty}(U)$ with $\psi^{\prime}\equiv 1$ on $\mathrm{supp}(\varphi^*\psi_1)$.
%For any $N$ linear symbols $\ell_{i}(x,hD_x) \in \Psi^1_{h}(\R^d)$, we define $$L_{i}=L_{\ell_{i}(x,hD_x),\psi_{i}}, \ \text{for any} \ i=1,\cdots,N.$$


As in Proposition \ref{inverseofelliptic}, we have
\begin{equation}\label{adfPHS}
\begin{aligned}
&\operatorname{ad}_{x_{i_1}} \cdots \operatorname{ad}_{x_{i_N}} \operatorname{ad}_{hD_{x_{i_{N+1}}}} \cdots \operatorname{ad}_{hD_{x_{i_{N+N^\prime}}}} (\varphi^{-1})^*\Phi \chi f(P)\chi\Phi^{-1}\varphi^*\\
&=(\varphi^{-1})^*\Phi \chi^{\prime} \big(\mathrm{ad}_{L_1}\cdots \mathrm{ad}_{L_{N+N'}} \chi f(P)\chi \big)\chi^{\prime}\Phi^{-1}\varphi^*\\
&= (\varphi^{-1})^*\Phi \chi^{\prime} \left(\frac{1}{\pi }\int_{\C}\overline{\partial}\widetilde{f}(z)\,\mathrm{ad}_{L_1}\cdots \mathrm{ad}_{L_{N+N'}} \chi (P-z)^{-1}\chi \,L(dz)\right)\chi^{\prime}\Phi^{-1}\varphi^*.
\end{aligned}
\end{equation}
By \eqref{adestiP-z} and \eqref{adfPHS}, we conclude that \eqref{pseudooflocal} holds for $f(P)$. Therefore, $f(P)\in \Psi^{0,sc}_{h}(M;F^\rho)$. The property $\sigma_{h}^{\rho}(f(P))=f(\sigma_{h}^{\rho}(P))$ is easily seen from the Helffer-Sj\"{o}strand formula.

\emph{Step 3. We prove that $\mathrm{WF}_{h}(f(P))\subset p_{0}^{-1}(\mathrm{supp}\,f)$ and $f(P)$ is compactly microlocalized.}  For any $(x_0,\xi_0)\notin p_{0}^{-1}(\mathrm{supp}\,f)$, there exist two open neighborhoods $U$ of $(x_0,\xi_0)$ and $V$ of $p_{0}^{-1}(\mathrm{supp}\,f)$ such that $U \cap V=\varnothing$. Let $T \in \Psi^{\mathrm{comp},sc}_{h}(M;F^\rho)$ with $\mathrm{WF}_{h}(T)\subset U$. Since 
\begin{equation*}
    \mathrm{WF}_{h}(T)\subset U\subset \mathrm{ell}_{h}(z-P), \qquad \forall z \in \mathrm{supp}\,\widetilde{f},
\end{equation*}
by the existence of elliptic parametrix \eqref{ellipticparametrix}, there exists $Q(z)\in \Psi^{\mathrm{comp},sc}_{h}(M;F^\rho)$ and $R(z)\in h^{\infty}\Psi^{-\infty,sc}(M;F^\rho)$ such that $T=(P-z)Q(z)+R(z)$ and $Q(z),\,R(z)$ are holomorphic for $z \in \mathrm{supp}\,\widetilde{f}$. Thus, we have
\begin{equation}
\begin{aligned}
f(P)T&=\frac{1}{\pi}\int_{\C}\overline{\partial}\widetilde{f}(z)(P-z)^{-1}T \, L(dz)\\
&=\frac{1}{\pi}\int_{\C}\overline{\partial}\widetilde{f}(z)Q(z)\,L(dz) + \frac{1}{\pi}\int_{\C}\overline{\partial}\widetilde{f}(z)(P-z)^{-1}R(z)\,L(dz)\\
&=\frac{1}{\pi}\int_{\C}\overline{\partial}\widetilde{f}(z)(P-z)^{-1}R(z)\, L(dz).
\end{aligned}
\end{equation}
By the almost analyticity of $\widetilde{f}$ and the estimate \eqref{sobolevnormofresolvent}, for any $z \in \mathrm{supp}\,\widetilde{f}$ and $N\in \mathbb{N}$, we have
\begin{equation*}
    \int_{\C}\|\overline{\partial}\widetilde{f}(z)(P-z)^{-1}\|_{H^N\to H^{N}}\,L(dz) \leq C_Nh^{-N},\qquad \|R(z)\|_{H^{-N}\to H^{N}}\leq C_Nh^{2N}.
\end{equation*}
We conclude that $f(P)T \in h^{\infty}\Psi^{-\infty,sc}(M;F^\rho)$. Therefore, $U\cap \mathrm{WF}_{h}(f(P))=\varnothing$ and $\mathrm{WF}_{h}(f(P))\subset p_0^{-1}(\mathrm{supp}\,f)$. From this and the fact that $|p_{0}(x,\xi)|\geq c\langle \xi \rangle_{x}^k-1$ for some $c>0$, we immediately obtain the compactness of $\mathrm{WF}_{h}(f(P))$.
\end{proof}

\end{document}
